考虑一维连续介质中的二分量狄拉克准粒子，速度参数 \(v>0\)，质量能隙尺度 \(\Delta>0\)，界面宽度 \(\xi>0\)，满足特殊可积条件 \(\Delta\xi=\hbar v\)。哈密顿量为
\[
H_0=-i\hbar v\sigma_x\partial_x+\Delta\tanh(x/\xi)\sigma_z,\qquad
\sigma_x=\begin{pmatrix}0&1\\1&0\end{pmatrix},\quad
\sigma_z=\begin{pmatrix}1&0\\0&-1\end{pmatrix}.
\]
取全空间 \(x\in\mathbb R\)，波函数平方可积且连续，\(\sigma_y=\begin{psmallmatrix}0&-i\\i&0\end{psmallmatrix}\)。在第三问中只考虑连续谱 \(|E|>\Delta\)，令 \(k=\sqrt{E^2/\Delta^2-1}>0\)，\(y=x/\xi\)。
(1) 证明 \(\{H_0,\sigma_y\}=0\)，直接解出归一化零能界面束缚态及其手征性；说明质量反号畴壁为何选取相反手征。
(2) 不借用已知谱表，严格平方 \(H_0\)，将两手征通道写成一对互伴的一阶算符乘积；证明在所给可积参数下除零模外没有其他离散束缚能级，并给出连续谱阈值。
(3) 构造非零能散射态的一个手征分量，利用互伴算符对自由平面波的作用，求透射振幅和反射概率。说明如何由该分量重构完整狄拉克自旋量，并检查概率流及 \(k\to0^+\)、\(k\to\infty\) 极限。
(4) 在全空间加入均匀手征破缺扰动 \(H_\delta=H_0+\delta\sigma_y\)，其中 \(\delta\in\mathbb R\) 具有能量量纲。求界面态的精确新能量、波函数及连续谱边缘，证明其仍在体能隙内；说明零能钉扎与界面局域化的区别。
